\section{Dataset Quality and Exploratory Data Analysis}
\label{sec:dataset_eda}

This section details data integrity, class distributions, temporal dynamics, feature representations, and the empirical validity of normality assumptions on the benchmark European credit card stream \cite{dalpozzolo2018credit}.

\subsection{Dataset Integrity and Chronological Partitioning}
The dataset comprises $284,807$ continuous transactions ($31$ numeric columns, $67.4\,\text{MB}$) spanning $48.0$ continuous hours in September 2013. A systematic sweep confirmed zero missing or null entries. We identified $1,081$ duplicate transactions ($0.38\%$), which are retained in the stream to reflect true production telemetry.

To prevent temporal data leakage, the stream is partitioned chronologically into 60\% training, 20\% validation, and 20\% test sets, as detailed in Table~\ref{tab:splits}. Notice that the fraud prevalence fluctuates naturally across time windows ($0.21\%$ in train, $0.10\%$ in validation, $0.13\%$ in test). Randomized $k$-fold cross-validation would artificially shuffle these non-stationary regimes, causing temporal snooping.

\begin{table}[h]
\centering
\caption{Temporal Split Statistics and Fraud Prevalence.}
\label{tab:splits}
\footnotesize
\begin{tabular}{lrrrr}
\hline
\textbf{Split Partition} & \textbf{Total Rows} & \textbf{Legit Rows} & \textbf{Fraud Rows} & \textbf{Base Rate (\%)} \\
\hline
Train ($[0, 120\,396\,\text{s}]$) & $170,884$ & $170,524$ & $360$ & $0.2107\%$ \\
Validation ($[120\,396, 145\,247\,\text{s}]$) & $56,961$ & $56,904$ & $57$ & $0.1001\%$ \\
Test ($[145\,248, 172\,792\,\text{s}]$) & $56,962$ & $56,887$ & $75$ & $0.1317\%$ \\
\hline
\textbf{Full Dataset} & $\mathbf{284,807}$ & $\mathbf{284,315}$ & $\mathbf{492}$ & $\mathbf{0.1727\%}$ \\
\hline
\end{tabular}
\end{table}

\subsection{Amount, Time, and Latent Dynamics}
\begin{itemize}
    \item \textbf{Transaction Amount:} Legitimate transactions average $\$88.29$ (median $\$22.00$, max $\$25,691$), while frauds average $\$122.21$ (median $\$9.25$, max $\$2,126$). Spending is severely right-skewed ($11.2\%$ IQR outliers in legitimate spending), requiring $\ln(1 + \text{Amount})$ transformation to stabilize variance.
    \item \textbf{Circadian Rhythms:} Aggregating elapsed `Time` into 1-hour bins demonstrates diurnal spending curves with activity dropping between 2--6 AM local time. Fraudulent activity remains comparatively steady overnight, exploiting off-peak staffing.
    \item \textbf{Discriminative Latent Features:} Two-sample Kolmogorov-Smirnov (KS) tests reveal strong separation along specific PCA dimensions: V14 ($\text{KS} = 0.8535$), V10 ($0.8320$), V12 ($0.7996$), V17 ($0.7870$), and V11 ($0.7854$). Off-diagonal correlations across legitimate training samples satisfy $|r| < 0.15$, confirming empirical orthogonality.
\end{itemize}

\subsection{Gaussianity Audit and Modeling Implications}
Parametric GMMs assume clusters conform to multivariate normality. We tested this assumption on legitimate training transactions across all 28 PCA dimensions:
\begin{itemize}
    \item \textbf{Normality Rejection:} Shapiro-Wilk tests rejected normality for all 28 features ($p < 0.05$).
    \item \textbf{Extreme Leptokurtosis:} While dimensions such as V11 and V13 exhibit near-Gaussian kurtosis ($\approx 0.15$), severe excess kurtosis dominates features like V2 ($\text{kurt} = 115.0$), V7 ($77.0$), and V8 ($261.7$). Only 9 of 28 features exhibit $|\text{kurt}| < 3.0$.
\end{itemize}
This empirical fat-tailed structure implies that single-component Gaussians ($K=1$) severely underestimate tail probability mass, establishing the theoretical need for multi-component mixtures ($K \ge 3$) or non-parametric baselines.

\subsection{Zero-Leakage Protocol Guarantee}
All preprocessing transformers (standardization, power scaling) are fit **strictly on legitimate training records** ($X_{\text{train}}[\text{Class}=0]$) and frozen for validation/test evaluation. Fraud labels are completely isolated during model estimation, and all triage thresholds are calibrated strictly on validation data.
